The ballot initiative pathway has produced more redistricting reform than any other mechanism. California, Colorado, and Michigan provide the most instructive examples: each achieved constitutional entrenchment of independent commissions, each used explicit partisan-balance requirements to ensure structural neutrality, and each established public engagement mandates that built legitimacy for the resulting maps.

None of the three states required algorithmic redistricting in their initial initiatives. Each created infrastructure---independent commissions, documented criteria, transparent processes---that is readily compatible with algorithmic requirements. The model language developed in Section~4 adds a \texttt{bisect} baseline requirement to California's existing CCRC framework without disturbing the Commission's authority or the VRA compliance hierarchy.

The political barriers to algorithmic ballot initiatives are real but manageable. The key framing insight is that the algorithm produces a \emph{baseline}, not a final map. Human commissioners retain authority to make adjustments for VRA compliance, community integrity, and other statutory criteria. The algorithm's role is to eliminate the manipulation space that exists between ``satisfies all stated criteria'' and ``maximizes partisan advantage''---and to do so transparently, in a way that any member of the public can verify.

For states considering redistricting reform initiatives, the following design checklist combines lessons from all three case studies:

\begin{enumerate}
  \item Amend the state constitution, not just statute.
  \item Create an independent commission with explicit partisan balance (not merely partisan neutrality).
  \item Adopt Michigan's sortition-based selection process where feasible.
  \item Enumerate criteria with explicit priority ordering, including anti-favoritism language from Colorado.
  \item Require public hearings at each stage with geographic distribution and remote access.
  \item Add the algorithmic baseline requirement from the model language in Section~4.
  \item Require documentation of departures from the baseline, with \textit{Gingles} justification for VRA adjustments.
  \item Mandate 12-year retention of all inputs, outputs, and manifests.
\end{enumerate}

The 2030 redistricting cycle begins in 2031. Initiative campaigns typically require 12--18 months from signature gathering to the ballot. States that wish to incorporate algorithmic baselines into their 2031 redistricting must begin the initiative process by 2028 at the latest. The technical infrastructure---the \texttt{bisect} platform---is already available. The design template is provided by the model language in Section~4. The remaining task is political organizing.
